\documentclass{article}
\input{./note-setup-leftsidebox}

\usepackage{setspace}
\usepackage{faktor}
\usepackage{etoolbox}
\usepackage{tikz-cd}
\usepackage{tensor}
\apptocmd{\normalsize}{%
  \setlength{\abovedisplayskip}{3pt}  
  \setlength{\belowdisplayskip}{3pt}  
  \setlength{\abovedisplayshortskip}{0pt} 
  \setlength{\belowdisplayshortskip}{3pt} 
  \setlength{\jot}{2pt}                
}{}{}
\DeclareMathOperator{\nil}{Nil}
\DeclareMathOperator{\spec}{Spec}
\DeclareMathOperator{\coim}{Coim}
\DeclareMathOperator{\coker}{Coker}
\DeclareMathOperator{\Hom}{Hom}
\DeclareMathOperator{\Mat}{Mat}
\DeclareMathOperator{\Nat}{Nat}
\setstretch{1.8}
\everymath{\displaystyle}
\title{Notes on GTM-73}
\date{\today}



\begin{document}

\maketitle
\newpage
\section{Introduction}
\begin{definition}[Cartesian Product]
    Let $\{A_i: i \in I\}$ be a family of sets indexed by a nonempty set $I$. The Cartesian product of the sets $A_{i}$ is the set of all functions $f:I\rightarrow \bigcup_{i\in I}A_{i}$ such that $f(i)\in A_{i}$ for all $i\in I$. It is denoted $\prod_{i\in I} A_{i}$.
\end{definition}

\textbf{The Axiom of Choice:} 
The product of a family of nonempty sets indexed by a nonempty set is nonempty.

There is another version of the Axiom of Choice.
\begin{theorem}
    Let $S$ be a set. A choice function for $S$ is a function $f$ from the set of all nonempty subsets of $S$ to $S$ such that $f(A)\in A$ for all $A\neq \emptyset $. Then the Axiom of Choice is equivalent to the statement thay every $S$ has a choice function.
\end{theorem}

\begin{proof}
    If the A.C. holds, let $\mathcal{P}(S)^*$ be the set of all nonempty sets of $S$. Then consider the product $\prod_{A\in \mathcal{P}(S)^*}A$. Since $A$ is nonempty, $\prod_{A\in \mathcal{P}(S)^*}A$ is nonempty. Take $f\in \prod_{A\in \mathcal{P}(S)^*}A$. According to the definition of Cartesian product, $f(A)\in A$, which also satisfies the definition of choice function. Hence every nonempty set $S$ has a choice function.
    
    Suppose $\{X_i\}$ is a family of nonempty sets, let $S=\bigcup_{i\in I}X_i$. Then there is a choice function for $S$ such that $f(A)\in A$ for all $A\neq \emptyset$ and $A\subseteq S$. Note that $X_i$ is also a nonempty set of $S$. Therefore $f(X_i)\in X_i$ for all $i\in I$. Then define $g$ as $g(i)=f(X_i)$. It is clearly that $g\in \prod_{i\in I}X_i$. Hence $\prod_{i\in I}X_i$ is nonempty.
\end{proof}


\begin{definition}
    A \textbf{partial ordering} on a set $A$ is a relation $R$ on $A\times A$ satisfies reflexive, transitive and antisymmetric. And $A$ with such a relation is called \textbf{partially ordered set}.
\end{definition}

Note that two gived elements of a partially ordered set need not to be comparable. If any two elements in a partially ordered set is comparable, we call the ordering a linear(or total or simple) ordering.

Here I do not give the specific definition of \textbf{maximal, upper bound} and \textbf{greatest element}. We should note that a maximal element is different from a greatest element. A nonempty subset $B$ of $A$ that is linearly ordered is called a \textbf{chain} in $A$.

\begin{lemma}[Zorn's lemma]
    If $A$ is a nonempty partially ordered set such that every chain in $A$ has an upper bound in $A$, then $A$ contains a maximal element.
\end{lemma}

\begin{definition}[Well ordered]
    If every nonempty subset of $A$ has a least element, then $A$ is said to be well ordered.
\end{definition}

Every well-ordered set is linearly ordered.


\newpage
\section{Ring}
\begin{definition}[Prime ideal]
    An ideal $P$ in a ring $R$ is said to be \textbf{prime} if $P\neq R$ and for any ideals $A,B$ in $R$, $AB \subseteq P$ implies $A\subseteq P$ or $B \subseteq P$.
\end{definition}

\begin{theorem}
    If $P$ is an ideal in a ring $R$ such that $P\neq R$ and for all $a,b \in R$,
        $ab \in P \Rightarrow a \in P $ or $ b \in P,$
    then $P$ is prime. Conversely if $P$ is prime and $R$ is commutative, then $P$ satisfies the condition above.
\end{theorem}

Consider $S$ to be the ring of all $n \times n$ matrices over a division ring $D$, $n \geq 2$. Then $S$ has no zero divisors and $0$ is a prime ideal which does not satisfies the condition of the theorem above. Thus the commutativity is necessary for the converse of the theorem.

\begin{theorem}
    In a nonzero ring $R$ with identity, maximal (left) ideals always exist. In fact every (left) ideal in $R$ except $R$ itself is contained in a maximal (left) ideal.
\end{theorem}
\begin{proof}
    The proof is application of Zorn's lemma. Note that set theoretic inclusion is a partial order.
\end{proof}

\begin{theorem}
    If $R$ is a commutative ring such that $R^2=R$ (in particular if $R$ has an identity), then every maximal ideal $M$ in $R$ is prime.
\end{theorem}

This theorem provides a sufficient condition for a maximal ideal to be a prime ideal. However, the maximal ideals seem "better" than prime ideals. Is this theorem meaningful?

\begin{theorem}
    Let $M$ be an ideal in a ring $R$ with identity $1_R\neq 0$.

    (1) If $M$ is maximal and $R$ is commutative, then the quotient ring $\faktor{R}{M}$ is a field.

    (2) If the quotient ring $\faktor{R}{M}$ is a division ring, then $M$ is a maximal ideal.
\end{theorem}
Note the conditions for applying the theorem. (1) is false when $R$ does not have an identity. If $M$ is maximal and $R$ is not commutative, then $\faktor{R}{M}$ need not to be a division ring.

\begin{definition}
    An element $e$ in a ring $R$ is said to be \textbf{idempotent} if $e^2=e$. An element of the center of the ring $R$ is said to be \textbf{central}. Idempotent elements $e_1, e_{2}, \ldots, e_{n}$ is said to be \textbf{orthogonal} if $e_{i}e_{j}=0$ for $i\neq j$.
\end{definition}

\begin{theorem}(Pierce decomposition)
    If $e$ is a central idempotent in a ring $R$ with identity, then $1_R-e$ is a central idempotent. $eR$ and $(1_R-e)R$ are ideals in $R$ such that $R=eR\times (1_R-e)R$.
\end{theorem}
This decomposition can be generalized to $R=\bigoplus_{i=1}^{n}\bigoplus_{j=1}^{n} e_{i}R e_{j} $ where $e_{1}, \ldots,e_{n}$ are idempotent and orthogonal and satisfy $e_{1}+e_{2}+ \cdots+e_{n}=1$.


\begin{definition}
    A nonempty subset $S$ of a ring $R$ is multiplicative provided that $a,b \in S$ implies $ab \in S$.
\end{definition}

\begin{theorem}
    Let $S$ be a multiplicative subset of a commutative ring $R$. The relation defined on the set $R\times S$ by 
    \begin{equation*}
        (r,s)\thicksim (r',s') \quad\iff\quad s_{1}(rs'-r's)=0\quad \text{for some}\quad s_{1}\in S
    \end{equation*}
    is an equivalence relation. Furthermore if $R$ has no zero divisors and $0 \notin S$, then
    \begin{equation*}
        (r,s)\thicksim (r',s') \quad\iff\quad rs'-r's=0.
    \end{equation*}
\end{theorem}
The $s_{1}$ in this relation aims to satisfy the transtive of equivalence relation.

Let $S^{-1}R$ be the set of equivalence classes of $R\times S$ under the equivalence relation of Theorem 2.6, then we call the ring $S^{-1}R$ (it is easy to verify that $S^{-1}R$ is a ring) the \textbf{ring of quotients}. If $S$ is the set of all nonzero elements of an integral domain $R$, then $S^{-1}R$ is a field called the \textbf{quotient field of the integral domain} $R$ (simply quotient field). 

\begin{theorem}
    Let $S$ be a multiplicative subset of a commutative ring $R$ and let $T$ be any commutative ring with identity. If $f:R \rightarrow T$ is a homomorphism of rings such that $f(s)$ is a unit in $T$ for all $s \in S$, then there exists a unique homomorphism of rings $\overline{f}:S^{-1}R \rightarrow T$ such that $\overline{f}\varphi_S = f$. The ring $S^{-1}R$ is completely determined (up to isomorphim) by this property.
\end{theorem}

\begin{example}
For instance, let $R$ be the integer ring $\mathbb{Z}$ and $S$ be the evens of $\mathbb{Z}$. Informally speaking, the ring of quotient $S^{-1}R=S^{-1}\mathbb{Z}$ is a ring such that every even is invertible in it. Take $X=\faktor{\mathbb{Z}[x]}{(2x-1)}$. Then in the ring $X$, $2n\equiv 1$ for any $n \in \mathbb{Z}$. Thus $X$ is also a ring such that every even in invertible. The theorem tells that there is a homomorphism between $S^{-1}\mathbb{Z}$ and $X$. Furthermore, the homomorphism is also an isomorphim, namely $S^{-1}\mathbb{Z} \simeq \faktor{\mathbb{Z}[x]}{(2x-1)}$.
\end{example}

\begin{theorem}
    Let $S$ be a multiplicative subset of a commutative ring $R$ with identity. Then there is a one-to-one correspondence between the set of prime ideals of $R$ which are disjoint from $S$ and the set of prime ideals of $S^{-1}R$, given by $P\mapsto S^{-1}R$.
\end{theorem}

\begin{definition}
    Let $R$ be a commutative ring with identity and $P$ is a prime ideal of $R$. Then $S=R-P$ is a multiplicative subset of $R$. The ring of quotients $S^{-1}R$ is called the \textbf{localization of $R$ at $P$} and is denoted as $R_P$. If $I$ is an ideal in $R$, then the ideal $S^{-1}I$ in $R_P$ is denoted as $I_P$.
\end{definition}

\begin{theorem}
    Let $P$ be a prime ideal of a commutative ring $R$ with identity. Then 

    (1) there is a one-to-one correspondence between the set of prime ideals of $R$ which are contained in $P$ and the set of prime ideals of $R_P$, given by $Q\mapsto Q_P$;

    (2) the ideal $P_P$ in $R_P$ is the unique maximal ideal of $R_P$.
\end{theorem}

\begin{definition}
    A \textbf{local ring} is a commutative ring with identity which has a unique maximal ideal.
\end{definition}

\begin{example}
    If $p$ is prime and $n\geqslant 1$, then $\mathbb{Z}_{p^n}$ is a local ring with maximal ideal $(p)$.
\end{example}
\begin{proof}
    A commutative ring with identity is a local ring if and only if all nonunits in the ring forms an ideal. Let $M$ be the set of all nonunits in $\mathbb{Z}_{p^n}$, then $M=\{m: p\mid m\}$ since all units in $\mathbb{Z}_{p^{n}}$ are relatively prime with $p$ and $p$ is prime. Thus $M=\mathbb{Z}_{p^n}p=(p)$. And $M=(p)$ is the unique maximal ideal.
\end{proof}

\begin{theorem}
    If $R$ is a commutative ring with identity then the following conditions are equivalent.

    (1) $R$ is a local ring;

    (2) all nonunits of $R$ are contained in some ideal $M\neq R$;

    (3) the nonunits of $R$ form an ideal.
\end{theorem}

\begin{theorem}
    In a commutative ring $R$ with identity the following conditions are equivalent:

    (1) $R$ has unique prime ideal;

    (2) every nonunit is nilpotent;

    (3) $R$ has minimal prime ideal which contains all zero divisors, and all nonunits of $R$ are zero divisors.
\end{theorem}
We need a lemma to prove this theorem.
\begin{lemma}
    The set of all nilpotents in $R$ equals to the intersection of all primes ideals of $R$, namely $\nil (R)=\bigcap_{P \in \spec(R)}P$.
\end{lemma}
The proof of this lemma and theorem can be found in the note on exercises.


\newpage
\section{Module}
\begin{theorem}
    Let $R$ be a ring and $A,A_{1},A_{2}, \ldots,A_{n}$ $R$-modules. Then $A\simeq \bigoplus_{i=1}^n A_{i}$ if and only if for each $i=1,2, \ldots,n$ there are $R$-modules homomorphism $\pi_i:A\to A_{i}$ and $\iota_i:A_{i}\to A$ such that

    (i) $\pi_i \iota_i = 1_{A_{i}}$ for $i=1,2, \ldots,n$;

    (ii) $\pi_{j} \iota_{i}=0$ for $i\neq j$;

    (iii) $\iota_{1}\pi_{1}+\iota_{2}\pi_{2}+ \cdots+\iota_{n}\pi_n=1_A$.
\end{theorem}
This actually provides a necessary and sufficient condition for the direct sum of $A_{1},A_{2}, \ldots,A_{n}$.

\begin{definition}
    A pair of module homomorphism, $A\xrightarrow{f}B\xrightarrow{g}C$, is said to be exact at $B$ provided that $\operatorname{Im}f=\operatorname{ker}g$. A finite sequence of module homomorphisms, $A_{0}\xrightarrow{f_{1}}A_{1}\xrightarrow{f_{2}}A_{2}\xrightarrow{f_{3}}\cdots\xrightarrow{f_{n-1}}A_{n-1}\xrightarrow{f_{n}}A_{n}$, is exact provided $\operatorname{Im}f_{i}=\operatorname{ker}f_{i+1} $ for $i=1,2, \ldots,n-1$. An infinite sequence of module homomorphisms, $\cdots\xrightarrow{f_{i-1}}A_{i-1}\xrightarrow{f_{i}}A_{i}\xrightarrow{f_{i+1}}A_{i+1}\xrightarrow{f_{i+2}}\cdots$ is exact provided $\operatorname{Im} f_{i}=\operatorname{ker} f_{i+1}$ for all $i \in \mathbb{Z}$.
\end{definition}

\begin{definition}
    If $f:A \to B$ is a module homomorphism, then $\faktor{A}{\operatorname{ker} f}$ is called the \textbf{coimage} of $f$, $\faktor{B}{\operatorname{Im} f}$ is called \textbf{cokernel} of $f$, denoted as $\coim f$ and $\coker f$ respectively.
\end{definition}

\begin{theorem}
    Let $R$ be a ring and $0\to A_{1}\xrightarrow{f}B\xrightarrow{g}A_{2}\to 0$ a short exact sequence of $R$-module homomorphisms. Then the short exact sequence is called \textbf{split} if it satisfies one of the following conditions:

    (1) there is an $R$-module homomorphism $h:A_{2}\to B$ such that $gh=\operatorname{id}_{A_{2}}$;

    (2) there is an $R$-module homomorphism $k:B\to A_{1}$ such that $fk=\operatorname{id}_{A_{1}}$;

    (3) the given sequence is equivalent to the short exact sequence $0\to A_{1}\xrightarrow{\iota_{1}}A_{1}\oplus A_{2}\xrightarrow{\pi_{2}}A_{2}\to 0$; in particular $B\simeq A_{1}\oplus A_{2}$.
\end{theorem}

\begin{theorem}
    Let $R$ be a ring with identity. $F$ is called \textbf{free $R$-module} if $F$ satisfies one of the following conditions:

    (1) $F$ has a nonempty basis;

    (2) $F$ is the internal direct sum of a family of cyclic $R$-modules, each of which is isomorphic as a left $R$-module to $R$;

    (3) $F$ is $R$-module isomorphic to a direct sum of copies of the left $R$-module R;

    (4) there exists a nonempyt set $X$ and a function $\iota:X \to F$ with the following property: given any unitary $R$-module $A$ and function $f:X\to A$, there exists a unique $R$-module homomorphism $\bar{f}:F\to A$ such that $\bar{f}\iota=f$. In other words, $F$ is a free object in the category of unitary $R$-modules.
\end{theorem}

\begin{corollary}
    Every (unitary) module $A$ over a ring $R$ (with identity) is the homomorphic image of a free $R$-module $F$. If $A$ is finitely generated, then $F$ may be chosen to be finitely generated.
\end{corollary}

\begin{example}
    Take $R=\mathbb{Z}$ and free $R$-module $F=R^2=\mathbb{Z}^2=\{(a,b): a,b \in \mathbb{Z}\}$. 
    
    Then $F$ has a basis $X=\{e_{1},e_{2}\}$, where $e_{1}=(1,0)$ and $e_{2}=(0,1)$. 
    
    Let $M_{1}=\mathbb{Z}e_{1}=\{(a,0):a \in \mathbb{Z}\}$, $M_{2}=\mathbb{Z}e_{2}=\{(0,b):b \in \mathbb{Z}\}$. Clearly $M_{1}$ and $M_{2}$ are cyclic $R$-modules and $M_{1} \simeq \mathbb{Z}$, $M_{2} \simeq \mathbb{Z}$. Since $M_{1}\cap M_{2}=\{0\}$ and $M_{1}+M_{2}=F$, then $F=M_{1}\oplus M_{2}$.

    Define $\iota: X \to \mathbb{Z}^2$ such that $\iota(x_{1})=(1,0)$ and $\iota(x_{2})=(0,1)$. Given an arbitrary $\mathbb{Z}-$module $A$ and a function $f:X \to A$, suppose $f(x_{1})=a_{1}$ and $f(x_{2})=a_{2}$. Define $\bar{f}:\mathbb{Z}^2\to A$ such that $\bar{f}(m,n)=ma_{1}+na_{2}$ where $m,n \in \mathbb{Z}$. It is easy to verfiy that $\bar{f}\iota=f$ and since the representation in terms of a basis is unique, the homomorphism $\bar{f}$ is uniquely determined.
\end{example}

Although $F$ may have a basis $X$ who has infinite elements, every element in $F$ can be represented as a finite sum of finite elements in $X$ since $F$ is isomorphic to a direct sum (coproduct) of copies of the modules.

Generally, we can define free module on any ring which may not have an identity. 
\begin{definition}
    Let $R$ be any ring and $X$ a nonempty set. An $R$-module $F$ is called a free module on $X$ if $F$ is a free object on $X$ in the category of all left $R$-modules.
\end{definition}

However, this defintion leads to a "bad" result.
\begin{example}
    Let $R$ be a ring with identity. Then $R$ is not a free module on any set in the category of all $R$-modules (as defined in definition 3.3).
\end{example}
\begin{proof}
    Let $A$ be a nonzero Abelian group with $R$-module structure $ra=0$ for any $r \in R$ and $a \in A$. Define $g:X\to A$ with $g(x_{0})=a_{0}\neq 0$. Then there exists a unique homomorphism $f:R\to A$ such that $f\iota=g$. For any $r$ in $R$, $f(r)=f(r\cdot 1_R)=rf(1_R)=0$ since $f(1_R)\in A$. Hence $g$ is a zero map, which contradicts to $g(x_{0})=a_{0}\neq 0$.
\end{proof}
If free module is defined as theorem 3.3 where the ring $R$ has an identity and $F$ is a unitary $R$-module, $R$ itself could be seen as a left free $R$-module with rank $1$.

\begin{definition}
    A module $P$ over a ring $R$ is said to be \textbf{projective} if given any diagram of $R$-modules homomorphisms
    $$\begin{tikzcd}
     & P\arrow[d,"f"] & \\
     A \arrow[r, "g"] & B \arrow[r]& 0\\ 
    \end{tikzcd}$$
    with bottom row exact, there exists an $R$-module homomorphism $h:P\to A$ such that the diagram 
    $$\begin{tikzcd}
                 & P \arrow[d, "f"] \arrow[ld, "h"'] &   \\
A \arrow[r, "g"] & B \arrow[r]                       & 0
\end{tikzcd}$$
is commutative (that is, $gh=f$).
\end{definition}

\begin{theorem}
    Every free module $F$ over a ring $R$ with identity is projective.
\end{theorem}
\begin{example}
    $R=\mathbb{Z}_{6}$, then $\mathbb{Z}_{3}$ and $\mathbb{Z}_{2}$are $\mathbb{Z}_{6}-$modules and it is easy to know that $\mathbb{Z}_{6}\simeq \mathbb{Z}_{2}\oplus \mathbb{Z}_{3}$. Hence both $\mathbb{Z}_{2}$ and $\mathbb{Z}_{3}$ are projective $\mathbb{Z}_{6}-$modules but are not free $\mathbb{Z}_{6}-$modules since there are zero divisors in $\mathbb{Z}_{6}$.
\end{example}

\begin{theorem}
    Let $R$ be ar ring. The following conditions on an $R$-module $P$ are equivalent.

    (1) $P$ is projective;

    (2) every short exact sequence $0\to A\xrightarrow{f}B\xrightarrow{g}P\to 0$ is split exact (hence $B\simeq A \oplus P$);

    (3) there is a free module $F$ and an $R$-module $K$ such that $F\simeq K\oplus P$.
\end{theorem}

\begin{definition}
    A module $J$ over a ring $R$ is said to be \textbf{injective} if given any diagram of $R$-module homomorphisms 
    $$\begin{tikzcd}
0 \arrow[r] & A \arrow[r, "g"] \arrow[d, "f"'] & B \\
            & J                                &  
\end{tikzcd}$$
with top row exact, there exists an $R$-module homomorphism $h:B\to J$ such that the diagram
$$\begin{tikzcd}
0 \arrow[r] & A \arrow[r, "g"] \arrow[d, "f"'] & B \arrow[ld, "h"] \\
            & J                                &                  
\end{tikzcd}$$
is commutaitve (that is, $hg=f$).
\end{definition}

\begin{theorem}
    Let $R$ be a ring. 

    (1) A direct sum of $R$-modules $\sum_{i\in I}P_i$ is projective iff each $P_i$ is projective;

    (2) A direct product of $R$-modules $\prod_{i\in I}J_{i}$ is injective iff each $J_{i}$ is injective.
\end{theorem}

\begin{theorem}[Baer Criterion]
    Let $R$ be a ring with identity. A unitary $R$-module $J$ is injective iff for every left ideal $L$ of $R$, any $R$-module homomorphism $L\to J$ may be extended to an $R$-module homomorphism $R\to J$.
\end{theorem}

\begin{theorem}
    Let $R$ be a ring with identity. The following conditions on a unitary $R$-module $J$ are equivalent.

    (1) $J$ is injective;

    (2) every short exact sequence $0\to J\xrightarrow{f}B\xrightarrow{g}C\to 0$ is split exact (hence $B\simeq J\oplus C$);

    (3) $J$ is a direct summand of any module $B$ of which it is a submodule.
\end{theorem}
\begin{proof}
    The detailed proof can be found in textbook.
\end{proof}

\begin{remark}
    (1) This theorem still holds without the identity. However, since both Baer's criterion and the existence of injective embeddings require the ring to be unital ring and the module to be unitary, the proof of theorem 3.9 is provided only for this case to ensure simplicity and utility.

    (2) Note that theorem 3.6(3) states that a projective module is a direct summand of some free module, whereas theorem 3.9(3) claims that an injective module is a direct summand of any module containing it. 
\end{remark}

\begin{example}
    $\mathbb{Q}$ is not a projective $\mathbb{Z}$-module.
\end{example}
\begin{proof}
    
\end{proof}
Is $\mathbb{Q}$ a projective $\mathbb{Q}$-module? The answer is yes since it is actually a vector space over a division ring, hence a free module.

Projective and injective modules are dual to each other. Is there, then, a concept dual to free modules? The following theorem addresses this question.

\begin{theorem}
    An $R$-module $F$ is co-free on a set $X$ if there exists a function $\iota:F\to X$ such that for any $R$-module $A$ and function $f:A\to X$, there exists a unique module homomorphism $\bar{f}:A\to F$ such that $\iota \bar{f}=f$. Then for any set $X$ with $\left\lvert X \right\rvert \geqslant 2$ no such $R$-module $F$ exists. If $\left\lvert X \right\rvert =1$, then $0$ is the only co-free module.
\end{theorem}
\begin{proof}
    Assume $F$ is co-free on a set $X$ with $\left\lvert X \right\rvert \geqslant 2$. Since $\bar{f}:A\to F$ is a homomorphism, then $\iota\bar{f}(0_A)=\iota(0_F)=f(0_A)$. Note that $f$ is an arbitrary function, then $\forall x \in X$, there exists $f:A\to X$ such that $f(0_A)=x$. But the function $\iota$ is fixed whence $\iota(0_F)$ is fixed. That is a contradiction. 
    
    If $\left\lvert X \right\rvert =1$, take $A=F$ then $\iota(k)=x, f(a)=x, \forall a \in A, k \in F$. Then for any homomorphism $\bar{f}:F\to F$, $\iota\bar{f}=f$ holds. In particular, when $\bar{f}=\operatorname{id}_{F}, 0$, $\iota\bar{f}=f$. Since $\bar{f}$ is unique, then $\operatorname{id}_{F}(k)=0(k)=0\Rightarrow k=0, \forall k \in F$. Hence $F$ is $0$ module.
\end{proof}

If $A,B$ are modules over a ring $R$, then $\Hom_{R}(A, B)$ is the set of all $R$-modules homomorphisms $f:A\to B$ and $\Hom_{R}(A, B)$ is an abelian group under addition and this addition is distributive with respect to composition of functions.

\begin{theorem}
    Let $A,B,C,D$ be modules over a ring $R$ and $\varphi:C\to A$ and $\psi:B\to D$ $R$-module homomorphism. Then the map $\theta:\Hom_{R}(A, B)\to \Hom_{R}(C, D)$ given by $f\mapsto \psi f \varphi$ is a homomorphism of abelian groups.
\end{theorem}
The map $\theta$ is usually denoted $\Hom_{}(\varphi, \psi)$ and called the homomorphism induced by $\varphi$ and $\psi$.

\begin{theorem}
    Let $R$ be a ring. $0\to A\xrightarrow{\varphi}B\xrightarrow{\psi}C$ is an exact sequence of $R$-modules iff for every $R$-module $D$
    $$0\to \Hom_{R}(D, A)\xrightarrow{\bar{\varphi}}\Hom_{R}(D, B)\xrightarrow{\bar{\psi}}\Hom_{R}(D, C)$$
    is an exact sequence of abelian groups.
\end{theorem}
\begin{theorem}
    Let $R$ be a ring. $A\xrightarrow{\theta}B\xrightarrow{\zeta}C\to 0$ is an exact sequence of $R$-modules iff for every $R$-module $D$
    $$0\to \Hom_{R}(C, D)\xrightarrow{\bar{\zeta}}\Hom_{R}(B, D)\xrightarrow{\bar{\theta}}\Hom_{R}(A, D)$$
    is an exact sequence of abelian groups.
\end{theorem}

By preceding results we say $\Hom_{R}(A, B)$ to be \textbf{left exact}. It is not true in general that a short sequence $0\to A\to B\to C\to 0$ induces a short exact sequence $0\to \Hom_{R}(D, A)\to \Hom_{R}(D, B)\to \Hom_{R}(D, C)\to 0$ or $0\to \Hom_{R}(C, D)\to \Hom_{R}(B, D)\to \Hom_{R}(A, D)\to 0$.

\begin{example}
    Consider $0\to \mathbb{Z}\xrightarrow{f}\mathbb{Z}\xrightarrow{g}\faktor{\mathbb{Z}}{2\mathbb{Z}}\to 0$ where $f(x)=2x$ and $g$ is canonical projection. 
    
    Take $D=\faktor{\mathbb{Z}}{2\mathbb{Z}}$, then $\Hom_{}(D, \faktor{\mathbb{Z}}{2\mathbb{Z}})\simeq \faktor{\mathbb{Z}}{2\mathbb{Z}}$ since homomorphism maps $0$ to $0$, and $\Hom_{}(D, \mathbb{Z})$ is zero since the order of the elements in $D$ is $2$ while there is no element in $\mathbb{Z}$ with order $2$. The induced sequence $0\to 0\xrightarrow{\bar{f}} 0\xrightarrow{\bar{g}}\faktor{\mathbb{Z}}{2\mathbb{Z}}\to 0$ is not an exact sequence since the map $\bar{g}$ is not an epimorphism.

    Take $D=\mathbb{Z}$, then $\Hom_{}(\mathbb{Z}, D)\simeq \mathbb{Z}$ since the homomorphism is determined by the identity in $\mathbb{Z}$, and $\Hom_{}(\faktor{\mathbb{Z}}{2\mathbb{Z}}, D)$ is zero. The induced sequence $0\to 0\xrightarrow{\bar{g}}\mathbb{Z}\xrightarrow{\bar{f}}\mathbb{Z}\to 0$ is not an exact sequence. In fact, $\bar{f}(\varphi)=\varphi f$ where $\varphi \in \Hom_{}(\mathbb{Z}, \mathbb{Z})$ and $f(x)=2x$, hence $\bar{f}$ could not be an epimorphism.
\end{example}

However, in some cases the exactness still holds in the induced sequence.


\begin{theorem}
    The following conditions on modules over a ring $R$ are equivalent.

    (1) $0\to A\xrightarrow{\varphi}B\xrightarrow{\psi}C\to 0$ is a split exact sequence of $R$-module;

    (2) $0\to \Hom_{R}(D, A)\xrightarrow{\bar{\varphi}} \Hom_{R}(D, C)\xrightarrow{\bar{\psi}}\Hom_{R}(D, C)\to 0$ is a split exact sequence of abelian groups for every $R$-module $D$;

    (3) $0\to \Hom_{R}(C, D)\xrightarrow{\bar{\psi}} \Hom_{R}(B, D)\xrightarrow{\bar{\varphi}}\Hom_{R}(A, D)\to 0$ is a split exact sequence of abelian groups for every $R$-module $D$.
\end{theorem}

\begin{corollary}
    The following conditions on a module $P$ over a ring $R$ are equivalent.

    (1) $P$ is projective;

    (2) if $\psi:B\to C$ is any $R$-module epimorphism, then $\bar{\psi}:\Hom_{R}(P, B)\to \Hom_{R}(P, C)$ is an epimorphism of abelian group;

    (3) if $0\to A\xrightarrow{\varphi}B\xrightarrow{\psi}C\to 0$ is any short exact sequence of $R$-modules, then $0\to \Hom_{R}(P, A)\xrightarrow{\bar{\varphi}} \Hom_{R}(P, C)\xrightarrow{\bar{\psi}}\Hom_{R}(P, C)\to 0$ is an exact sequence of abelian groups.
\end{corollary}

\begin{corollary}
    The following conditions on a module $J$ over a ring $R$ are equivalent.

    (1) $J$ is injective;

    (2) if $\theta:A\to B$ is any $R$-module monomorphism, then $\bar{\theta}:\Hom_{R}(B, J)\to \Hom_{R}(A, J)$ is an epimorphism of abelian group;

    (3) if $0\to A\xrightarrow{\theta}B\xrightarrow{\zeta}C\to 0$ is any short exact sequence of $R$-modules, then $0\to \Hom_{R}(C, J)\xrightarrow{\bar{\zeta}} \Hom_{R}(B, J)\xrightarrow{\bar{\theta}}\Hom_{R}(A, J)\to 0$ is an exact sequence of abelian groups.
\end{corollary}


Intuitively, $\Hom$ functors are only left-exact and lose surjectivity at the right end. A projective module "repairs" this for the covariant functor via its lifting property, while an injective module "repairs" this for the contravariant functor (which reverses arrows) via its extending property.

\begin{theorem}
    Let $A,B,\{A_{i}:i \in I\}$ and $\{B_{j}:j \in J\}$ be modules over a ring $R$. Then there are isomorphims of abelian groups:

    (1) $\Hom_{R}(\bigoplus_{i \in I}A_{i}, B)\simeq \prod_{i \in I} \Hom_{R}(A_{i}, B)  $;

    (2) $\Hom_{R}(A, \prod_{j \in J}B_{j}  )\simeq \prod_{j \in J}\Hom_{R}(A, B_{j})  $.
\end{theorem}

\begin{example}
    Let $R^n$ and $R^m$ be free modules of rank $n$ and $m$ respectively, then $R^n = \bigoplus_{j=1}^n R$ and $R^m = \prod_{i=1}^m R$ since products and coproducts are isomorphic under finite modules. By the theorem above, $$\Hom_R\left( \bigoplus_{j=1}^n R, R^m \right) \simeq \prod_{j=1}^n \Hom_R(R, R^m)\simeq \prod_{j=1}^{n} \prod_{i=1}^{m} \Hom_{R}(R, R) \simeq  \prod_{j=1}^{n} \left( \prod_{i=1}^{m} R\right). $$ This corresponds to the matrix of $m \times n$ on $R$.
\end{example}

Since $\Hom_{R}(A, B)$ forms an abelian group not a module, we need consider possible module structures on it.

Let $R$ and $S$ be rings. An abelian group $A$ is said to be an $R$-$S$ \textbf{bimodule} provided that $A$ is both a left $R$-module and a right $S$-module and $r(as)=(ra)s$ for all $a \in A, r\in R, s \in S$, denoted as $\tensor[_R]{A}{_S}$.

\begin{theorem}
    Let $R$ and $S$ be rings and let $\tensor[_R]{A}{}, \tensor[_R]{B}{_S}, \tensor[_R]{C}{_S}, \tensor[_R]{D}{}$ be (bi)modules as indicated.

    (i) $\Hom_{R}(A, B)$ is a right $S$-module, with the action of $S$ given by $(fs)(a)=(f(a))s,s \in S, a \in A, f \in \Hom_{R}(A, B)$.

    (ii) If $\varphi:A\to A'$ is a homomorphism of left $R$-modules, then the induced map $\bar{\varphi}:\Hom_{R}(A', B)\to \Hom_{R}(A, B)$ is a homomorphism of right $S$-modules.

    (iii) $\Hom_{R}(C, D)$  is a left $S$-module, with the action of $S$ given by $(sg)(c)=g(cs), s \in S, c \in C, g \in \Hom_{R}(C, D)$.

    (iv) If $\psi:D\to D'$ is a homomorphism of left $R$-module, then $\bar{\psi}:\Hom_{R}(C, D)\to \Hom_{R}(C, D')$ is a homomorphism of left $S$-modules.
\end{theorem}

\begin{theorem}
    If $A$ is a unitary left module over a ring $R$ with identity, then there is an isomorphism of left $R$-modules $A \simeq \Hom_{R}(R, A)$.
\end{theorem}

Through this theorem, we can easily prove that a ring $R$ with identity is projective as a $R$-module.

\begin{corollary}
    A ring $R$ with identity is projective as a left $R$-module.
\end{corollary}
\begin{proof}
    Given a short exact sequence $0\to A\xrightarrow{\varphi}B\xrightarrow{\psi}C\to 0$ of $R$-modules, if $0\to \Hom_{R}(R, A)\xrightarrow{\bar{\varphi}}\Hom_{R}(R, B)\xrightarrow{\bar{\psi}}\Hom_{R}(R, C)\to 0$ is also an exact sequence then $R$ is projective by Corollary 3.15.
    
    By Theorem 3.19, $\Hom_{R}(R, X)\simeq X, X=A,B,C$. Clearly $\Hom_{R}(R, -)$ keeps the exactness of the sequence (the diagram is commutative). Thus $R$ is projective. 
\end{proof}

\begin{definition}
    Let $A$ be a left module over a ring $R$, then $\Hom_{R}(A, R)$ is a right $R$-module and is called the \textbf{dual} module of $A$, denoted as $A^*$.
\end{definition}

\begin{theorem}
    Let $A,B,C$ be left modules over a ring $R$.

    (i) If $\varphi:A\to C$ is a homomorphism of left $R$-modules, then the induced map $\bar{\varphi}:C^*=\Hom_{R}(C, R)\to \Hom_{R}(A, R)=A^*$ is a homomorphism of right $R$-module.

    (ii) There is an $R$-module isomorphism $(A\oplus C)^*\simeq A^*\oplus C^*$.

    (iii) If $R$ is a division ring and $0\to A\xrightarrow{\theta}B\xrightarrow{\zeta}C\to 0$ is a short exact sequence of left vector spaces, then $0\to C^*\xrightarrow{\bar{\zeta}}B^*\xrightarrow{\bar{\theta}}A^*\to 0$ is a short exact sequence of right vector spaces.
\end{theorem}
The map $\bar{\varphi}$ is called the \textbf{dual map} of $\varphi$.

\begin{theorem}
    Let $F$ be a free module over a ring $R$ with identity. Let $X$ be a basis of $F$ and for each $x \in X$ let $f_x:F\to R$ be given by $f_x(y)=\delta_{xy}$. Then 

    (i) $\{f_x:x \in X\}$ is a linearly independent subset of $F^*$ of cardinality $\left\lvert X \right\rvert $;

    (ii) if $X$ is finite, then $F^*$ is a free right $R$-module with basis $\{f_x:x \in X\}$.
\end{theorem}

The process of forming duals may be repeated. If $A$ is a left $R$-module, then $A^*$is a right $R$-module and $A^{**}=(A^*)^*=\Hom_{R}(\Hom_{R}(A, R), R)$ is a left $R$-module. $A^{**}$ is called the double dual of $A$.

\begin{theorem}
    Let $A$ be a left module over a ring $R$.

    (i) There is an $R$-module homomorphism $\theta:A\to A^{**}$.

    (ii) If $R$ has an identity and $A$ is free, then $\theta$ is a monomorphism.

    (iii) If $R$ has an identity and $A$ is free with a finit basis, then $\theta$ is an isomorphism.
\end{theorem}
A module $A$ such that $\theta:A \to A^{**}$ is an isomorphism is said to be \textbf{reflexive}.
\begin{example}
    The theorem above without the condition that $A$ has a finite basis does not hold. Let $F=\bigoplus_{x \in X}\mathbb{Z}x$ be a free $\mathbb{Z}$-module with an infinite basis $X$. Then $\{f_x:x \in X\}$ does not form a basis of $F^*$.
\end{example}
\begin{proof}
    Note that $F^*=\Hom_{\mathbb{Z}}(F, \mathbb{Z})=\Hom_{\mathbb{Z}}(\bigoplus_{x \in X} \mathbb{Z}x, \mathbb{Z}) \simeq \prod_{x \in X}\Hom_{\mathbb{Z}}(\mathbb{Z}x, \mathbb{Z}) \simeq \prod_{x \in X} \mathbb{Z}x  $ which is a direct product, not a direct sum.
\end{proof}

\begin{definition}
    Let $A_R$ and $\tensor[_R]{B}{}$ be modules over a ring $R$ and $C$ be an additive abelian group, then a \textbf{middle linear map} from $A\times B$ to $C$ is a function $f:A\times B\to C$ such that for all $a,a_{i}\in A,b,b_{i}\in B$ and $r \in R$:
    \begin{equation*}
        \begin{aligned}
        f(a_{1}+a_{2},b) &=f(a_{1},b) + f(a_{2}, b) \\
        f(a,b_{1}+b_{2}) &=f(a,b_{1}) + f(a,b_{2}) \\
        f(ar,b) &= f(a,rb).
        \end{aligned}
    \end{equation*}
\end{definition}

\begin{definition}
    Let $A$ be a right module and $B$ a left module over a ring $R$. Let $F$ be the free abelian group on the set $A\times B$. Let $K$ be the subgroup of $F$ generated by all elements of the following forms (for all $a,a' \in A$; $b,b' \in B$; $r \in R$):

    (i) $(a+a',b)-(a,b)-(a',b)$;

    (ii) $(a,b+b')-(a,b)-(a,b')$;

    (iii) $(ar,b)-(a,rb)$.

    The quotient group $\faktor{F}{K}$ is called the \textbf{tensor product} of $A$ and $B$; it is denoted $A\otimes_R B$ (or simply $A\otimes B$ if $R=\mathbb{Z}$). The coset $(a,b)+K$ of the element $(a,b)$ in $F$ is denoted $a\otimes b$; the coset of $(0,0)$ is denoted $0$.
\end{definition}

The element of $A\otimes_R B$ usually is of the form $\sum_{i=1}^{r}n_{i}(a_{i}\otimes b_{i})$ and note that the tensor product is the coset of $\faktor{F}{K}$ hence may has different representatives. Also by defintion, for all $a,a_{i} \in A$, $b,b_{i} \in B$, and $r \in R$, there follows:
\begin{equation*}
    \begin{aligned}
    (a_{1}+a_{2})\otimes b &= a_{1}\otimes b + a_{2}\otimes b; \\
    a\otimes (b_{1}+b_{2}) &= a\otimes b_{1} + a\otimes b_{2}; \\
    ar\otimes b &= a\otimes rb.
    \end{aligned}
\end{equation*} 

Given modules $A_R$ and $\tensor[_R]{B}{}$ over a ring $R$, it is easy to verify that the map $i:A\times B\to A\otimes_R B$ given by $(a,b)\mapsto a\otimes b$ is a middle linear map. The map $i$ is called the \textbf{canonical middle linear map}.

\begin{theorem}
    Let $A_R$ and $\tensor[_R]{B}{}$ be modules over a ring $R$, and let $C$ be an abelian group. If $g:A\times B\to C$ is a middle linear map, then there exists a unique group homomorphism $\bar{g}:A\otimes_R B\to C$ such that $\bar{g}i=g$, where $i:A\times B \to A\otimes_R B$ is the canonical middle linear map. $A\otimes_R B$ is uniquely determined up to isomorphism by this property. In other words $i:A\times B\to A\otimes_R B$ is universal in the category $\mathfrak{M} (A,B)$ of all middle linear map on $A\times B$.
\end{theorem}

\begin{theorem}
    If $A\xrightarrow{f}B\xrightarrow{g}C\to 0$ is an exact sequence of left moduels over a ring $R$ and $D$ is a right $R$-module, then 
    \begin{equation*}
        D\otimes_R A\xrightarrow{1_D\otimes f}D\otimes_R B\xrightarrow{1_D\otimes g}D\otimes_R C\to C
    \end{equation*}
    is an exact sequence of abelian groups. An analogous statement holds for an exact sequence in the first varible.
\end{theorem}

The tensor product preserves surjections. If $h:A_R\to A_R'$ and $k:\tensor[_R]{B}{}\to\tensor[_R]{B}{}'$ are modules epimorphisms, then by Theorem 3.25, $h\otimes k: A\otimes_R B\to A'\otimes_R B'$ is an epimorphism since $h\otimes k=(1_{A'}\otimes k)(h\otimes 1_B)$. However, if $h$ and $k$ are monomorphisms, $1_{A'}\otimes k$ and $h\otimes 1_B$ need not to be monomorphisms.
\begin{example}
    The usual injection $\alpha:\mathbb{Z}_2\to \mathbb{Z}_4$ is a monomorphism of abelian group. But the map $1\otimes \alpha: \mathbb{Z}_2\otimes \mathbb{Z}_2\to \mathbb{Z}_2\otimes \mathbb{Z}_4$ is the zero map.
\end{example}
\begin{proof}
    Since $\alpha$ is a homomorphism, then $\alpha(x)=2x,x=\bar{0},\bar{1}$. Thus $(1\otimes \alpha)(a\otimes b)=a\otimes 2b$ where $a,b=\bar{0},\bar{1}$. Since $0=2\cdot 0$, then $1\otimes \alpha$ is the zero map.
\end{proof}

\begin{theorem}
    Let $0\to A\xrightarrow{f}B\xrightarrow{g}C\to 0$ is a short exact sequence of left $R$-module and $D$ a right $R$-module. Then $0\to D\otimes_R A\xrightarrow{1_D \otimes f}D\otimes_R B\xrightarrow{1_D \otimes g}D \otimes_R C\to 0$ is a short exact sequence of abelian groups if one of the following conditions holds:
    
    (i) $0\to A\xrightarrow{f}B\xrightarrow{g}C\to 0$ is split exact.

    (ii) $R$ has an identity and $D$ is a free right $R$-module.

    (iii) $R$ has an identity and $D$ is a projective unitary right $R$-module.
\end{theorem}
\begin{proof}
    
\end{proof}

\begin{theorem}
    Let $R$ and $S$ be rings and $\tensor[_S]{A}{_R},\tensor[_R]{B}{},C_R,\tensor[_R]{D}{_S}$ (bi)modules as indicated.

    (i) $A\otimes_R B$ is a left $S$-module such that $s(a\otimes b)=sa\otimes b$ for all $s \in S, a \in A, b \in B$.

    (ii) If $f:A\to A'$ is a homomorphism of $S-R$ bimodules and $g:B\to B'$ is an $R$-module homomorphism, then the induced map $f\otimes g:A\otimes_R B\to A'\otimes_R B'$ is a homomorphism of left $S$-modules.

    (iii) $C\otimes_R D$ is a right $S$-module such that $(c\otimes d)s=c\otimes ds$ for all $c \in C, d \in D, s \in S$.

    (iv) If $h:C\to C'$ is an $R$-module homomorphism and $k:D\to D'$ a homomorphism of $R-S$ bimodules, then the induce map $h\otimes k:C\otimes_R D\to C'\otimes_R D'$ is a homomorphism of right $S$-module.
\end{theorem}

If $R$ is a commutative ring, then the tensor product of $R$-module may be characterized by a useful variation. Let $A,B,C$ be modules over a commutative ring $R$. A \textbf{bilinear map} from $A\times B$ to $C$ is a function $f:A\times B\to C$ such that for all $a,a_{i} \in A, b,b_{i} \in B$, and $r \in R$:
\begin{equation*}
    \begin{aligned}
    f(a_{1}+a_{2},b) &= f(a_{1},b) + f(a_{2},b) \\
    f(a,b_{1}+b_{2}) &= f(a,b_{1}) + f(a,b_{2}) \\
    f(ra,b) &= f(a,rb) =rf(a,b).
    \end{aligned}
\end{equation*}

If $A$ and $B$ are modules over a commutative ring $R$, then so is $A\otimes_R B$ and the canonical middle linear map $i:A\times B\to A\otimes_R B$ is easily seen to be bilinear. In this context $i$ is called the \textbf{canonical bilinear map}.

\begin{theorem}
    If $A,B,C$ are modules over a commutative ring $R$ and $g:A\times B\to C$ is a bilinear map, then there is a unique $R$-module homomorphism $\bar{g}:A\otimes_R B\to C$ such that $\bar{g}i=g$, where $i:A\times B\to A\otimes_R B$ is the canoncial bilinear map. The module $A \otimes_R B$ is uniquely determined up to isomorphism by this property. 
\end{theorem}

\begin{theorem}
    If $R$ is a ring with identity and $A_R, \tensor[_R]{B}{}$ are unitary $R$-modules, then there are $R$-module isomophisms
    \begin{equation*}
        A \otimes_R R \simeq A \quad \text{and} \quad R \otimes_R B \simeq B.
    \end{equation*}
\end{theorem}

\begin{theorem}
    If $R$ and $S$ are rings and $A_R,\tensor[_R]{B}{_S},\tensor[_S]{C}{}$ are (bi)modules, then there is an isomorphism 
    \begin{equation*}
        (A \otimes_R B) \otimes_S C \simeq A \otimes_R (B \otimes_S C).
    \end{equation*}
\end{theorem}
\begin{proof}
    Fix $c \in C$ and define $f_c:A\times B\to A \otimes_R (B \otimes_S C)$ such that $f_c(a,b)=a \otimes (b \otimes c)$. Then $f_c(ar, b) = (ar) \otimes (b \otimes c) = a \otimes r(b \otimes c) = a \otimes (rb \otimes c) = f_c(a, rb)$. Thus there exists a unique group homomorphism $\varphi_c:A \otimes_R B \to A \otimes_R (B \otimes_S C)$. 
    
    Define $g:(A \otimes_R B)\times C\to A \otimes_R (B \otimes_S C)$ such that $g(x,c)=\varphi_c(x)$ where $x \in A \otimes_R B$. Then for $x=a \otimes b$, $g(xs,c)=\varphi_c(xs)=\varphi_c((a \otimes b)s)=\varphi_c(a \otimes bs)=a \otimes (bs \otimes c)=a \otimes (b \otimes sc)=\varphi_{sc}(x)=g(x,sc)$. Hence there exists a unique group homomorhpism $\phi:(A \otimes_R B) \otimes_S C\to A \otimes_R (B \otimes_S C)$ such that $\phi((a \otimes b)\otimes c)=a \otimes (b \otimes c)$.

    Symmetrically we can construct the inverse map $\psi:A \otimes_R (B \otimes_S C)\to (A \otimes_R B) \otimes_S C$ such that $\psi(a \otimes (b \otimes c))=(a \otimes b) \otimes c$. Note that $\phi \psi =\operatorname{id}_{}$ and $\psi \phi =\operatorname{id}_{}$. Hence $(A \otimes_R B) \otimes_S C \simeq A \otimes_R (B \otimes_S C)$.
\end{proof}

\begin{theorem}
    Let $R$ be a ring, $A$ and $\{A_{i}:i \in I\}$ right $R$-modules, $B$ and $\{B_{j}:j \in J\}$ left $R$-modules. Then there are group isomorphism.

    $$\begin{aligned}
    (\bigoplus_{i \in I} A_{i}) \otimes_R B &\simeq \bigoplus_{i \in I}(A_{i}\otimes_R B); \\
    A\otimes_R (\bigoplus_{j \in J} B_{j}) &\simeq \bigoplus_{j \in J} (A\otimes_R B_{j}).
    \end{aligned}$$
\end{theorem}
\begin{proof}
    $$\begin{tikzcd}
(\bigoplus_{i\in I}A_i)\times B \arrow[rr, "f"] \arrow[d, "i"'] &  & \bigoplus_{i\in I}(A_i\otimes_R B) \\
(\bigoplus_{i\in I}A_i)\otimes_R B \arrow[rru, "\varphi"']      &  &                                   
\end{tikzcd}$$
$$    \begin{tikzcd}
A_i\times B \arrow[rr, "\psi_i"] \arrow[d, "i_2"']              &  & (\bigoplus_{i\in I} A_i)\otimes_R B \arrow[d, equal] \\
A_i\otimes_R B \arrow[d, "\iota_i"'] \arrow[rr, "\bar{\psi_i}"] &  & (\bigoplus_{i\in I} A_i)\otimes_R B                                \\
\bigoplus_{i\in I}(A_i\otimes_R B) \arrow[rru, "\psi"']         &  &                                                                   
\end{tikzcd}$$
$\varphi \psi=\operatorname{id}_{}, \psi \varphi=\operatorname{id}_{}$.
%Step 1: Construction of $\Phi$Define a map $\varphi: (\bigoplus_{i \in I} A_i) \times B \to \bigoplus_{i \in I} (A_i \otimes_R B)$ by:$$\varphi((a_i)_{i \in I}, b) = (a_i \otimes b)_{i \in I}$$Since $(a_i)_{i \in I}$ has only finitely many non-zero components, the sequence $(a_i \otimes b)_{i \in I}$ also has finitely many non-zero components, making it a well-defined element of the direct sum. It is straightforward to verify that $\varphi$ is $R$-middle linear. By the universal property of the tensor product, there exists a unique homomorphism $\Phi: (\bigoplus_{i \in I} A_i) \otimes_R B \to \bigoplus_{i \in I} (A_i \otimes_R B)$ such that:$$\Phi((a_i)_{i \in I} \otimes b) = (a_i \otimes b)_{i \in I}$$Step 2: Construction of $\Psi$For each $j \in I$, let $\iota_j: A_j \to \bigoplus_{i \in I} A_i$ be the canonical injection. The assignment $(a_j, b) \mapsto \iota_j(a_j) \otimes b$ is an $R$-middle linear map from $A_j \times B$ to $(\bigoplus_{i \in I} A_i) \otimes_R B$. By the universal property of the tensor product, this induces a unique homomorphism $f_j: A_j \otimes_R B \to (\bigoplus_{i \in I} A_i) \otimes_R B$ such that:$$f_j(a_j \otimes b) = \iota_j(a_j) \otimes b$$By the universal property of the direct sum (coproduct), the family of homomorphisms $\{f_j\}_{j \in I}$ induces a unique homomorphism $\Psi: \bigoplus_{i \in I} (A_i \otimes_R B) \to (\bigoplus_{i \in I} A_i) \otimes_R B$ whose restriction to the $j$-th summand is $f_j$. Explicitly, on a generator belonging to the $j$-th summand, we have:$$\Psi(a_j \otimes b) = \iota_j(a_j) \otimes b$$Step 3: Verification of IsomorphismIt suffices to check that $\Psi \circ \Phi$ and $\Phi \circ \Psi$ are identity maps on their respective sets of generators.For any generator $(a_i)_{i \in I} \otimes b \in (\bigoplus_{i \in I} A_i) \otimes_R B$:$$\Psi(\Phi((a_i)_{i \in I} \otimes b)) = \Psi((a_i \otimes b)_{i \in I})$$Since the element $(a_i \otimes b)_{i \in I}$ has finitely many non-zero terms, we can expand it as a finite sum. By the additivity of $\Psi$:$$\Psi((a_i \otimes b)_{i \in I}) = \sum_{j} \Psi(a_j \otimes b) = \sum_{j} (\iota_j(a_j) \otimes b)$$By the additivity of the tensor product in the first variable, we can factor out $b$:$$\sum_{j} (\iota_j(a_j) \otimes b) = \left(\sum_{j} \iota_j(a_j)\right) \otimes b = (a_i)_{i \in I} \otimes b$$Hence, $\Psi \circ \Phi = \text{id}$.Conversely, for any generator $a_j \otimes b$ in the $j$-th summand of $\bigoplus_{i \in I} (A_i \otimes_R B)$:$$\Phi(\Psi(a_j \otimes b)) = \Phi(\iota_j(a_j) \otimes b)$$By the definition of $\Phi$, evaluating this yields a sequence with $a_j \otimes b$ in the $j$-th position and $0$ elsewhere, which is precisely the generator we started with. Hence, $\Phi \circ \Psi = \text{id}$.Therefore, $\Phi$ is an isomorphism. 
\end{proof}

\begin{theorem}[Adjoint Associativity]
    Let $R $ and $S$ be rings and $A_R,\tensor[_R]{B}{_S},C_S$ (bi)modules. Then there is an isomorphism of abelian groups
    \begin{equation*}
        \alpha: \Hom_{S}(A \otimes_R B, C) \simeq \Hom_{R}(A, \Hom_{S}(B, C)),
    \end{equation*}
    defined for each $f:A \otimes_R B\to C$ by $$[(\alpha f)(a)](b)=f(a \otimes b).$$
\end{theorem}
\begin{proof}
    For all $a, a_{i} \in A, b, b_{i} \in B, f,f_{i} \in \Hom_{S}(A\otimes_R B, C),r \in R, s \in S$, we need to check that $(\alpha(f))(a)$ is a $S$-module homomorphism, $\alpha(f)$ is an $R$-module homomorphism and $\alpha$ is a group homomorphism in turn.
    \begin{equation*}
        \begin{aligned}
        ((\alpha f)(a))(b_{1}+b_{2}) &=f(a\otimes(b_{1}+b_{2}))=f(a\otimes b_{1}+a\otimes b_{2}) \\
        & =f(a\otimes b_{1})+f(a\otimes b_{2})=((\alpha f)(a))(b_{1})+((\alpha f)(a))(b_{2}), \\
        ((\alpha f)(a))(bs) &= f(a\otimes bs)=f((a \otimes b )s)=f(a \otimes b)s=((\alpha f)(a))(b)s.
        \end{aligned}
    \end{equation*}
    Hence $(\alpha f)(a)$ is a $S$-module homomorphism.
    $$        ((\alpha f)(a_{1}+a_{2}))(b) = f((a_{1}+a_{2})\otimes b)=f(a_{1} \otimes b)+f(a_{2} \otimes b)=((\alpha f)(a_{1}))(b)+((\alpha f)(a_{2}))(b),$$ 
    $$        ((\alpha f)(ar))(b) = f(ar \otimes b)=f(a \otimes rb)=((\alpha f)(a))(rb)=[(\alpha f)(a)\cdot r](b).$$
    Hence $\alpha f$ is an $R$-module homomorphism.
    \begin{equation*}
        (\alpha (f_{1}+f_{2})(a))(b)=(f_{1}+f_{2})(a \otimes b)=f_{1}(a \otimes b) + f_{2}(a \otimes b)=((\alpha f_{1})(a))(b)+ ((\alpha f_{2})(a))(b),
    \end{equation*}
    thus $\alpha$ is a group homomorphism.

    Now we need to find its inverse map $\beta:\Hom_{R}(A, \Hom_{S}(B, C)) \to\Hom_{S}(A \otimes_R B, C) $. Note that for all $g \in \Hom_{R}(A, \Hom_{S}(B, C))$, there is $\beta g \in \Hom_{S}(A\otimes_R B, C)$ and $\beta g$ is a map from a tensor product to a set. Thus we first define $\psi:A\times B \to C$ such that $\psi (a,b)=(g(a))(b)$. Since $A\times B$ is Cartesian product and $g \in \Hom_{S}(A\otimes_R B, C)$, then easily $\psi$ is well-defined.
    $$\psi(a_{1}+a_{2},b)=(g(a_{1}+a_{2}))(b)=(g(a_{1}))(b)+(g(a_{2}))(b)=\psi(a_{1},b)+\psi(a_{2},b),$$
    $$\psi(a,b_{1}+b_{2})=(g(a))(b_{1}+b_{2})=(g(a))(b_{1})+(g(a))(b_{2})=\psi(a,b_{1})+\psi(a,b_{2}),$$
    $$\psi(ar,b)=(g(ar))(b)=(g(a)r)(b)=(g(a))(rb)=\psi(a,rb).$$
    Hence $\psi$ is a middle linear map and there exists a unique group homomorphism $\beta g:A \otimes_R B\to C$ such that $(\beta g)(a \otimes b)=(g(a))(b)$.
    $$(\beta g)((a\otimes b)s)=(\beta g)(a \otimes bs)=(g(a))(bs)=(g(a))(b)\cdot s=(\beta g)(a \otimes b)s,$$
    thus $\beta g$ is also a $S$-module homomorphism.
    
    For all $g_{1},g_{2} \in \Hom_{S}(A\otimes_R B, C)$, there is 
    $$(\beta (g_{1}+g_{2}))(a \otimes b)=((g_{1}+g_{2})(a))(b)=(g_{1}(a))(b)+(g_{2}(a))(b)=(\beta g_{1})(a_{1} \otimes b)+(\beta g_{2})(a_{2} \otimes b).$$
    Hence $\beta $ is a group homomorphism.
    Then $[(\alpha (\beta g))(a)](b)=(\beta g)(a \otimes b)=(g(a))(b)$ implies $\alpha (\beta g)=(\alpha \beta)g=g$ and $(\beta (\alpha f))(a \otimes b)=((\alpha f)(a))(b)=f(a \otimes b)$ imples $\beta (\alpha f)=(\beta \alpha )f =f$. Hence $\alpha$ is a group isomorphism and $\Hom_{S}(A \otimes_R B, C) \simeq \Hom_{R}(A, \Hom_{S}(B, C))$.
\end{proof}

\begin{theorem}
    Let $A,B,C$ be modules over a commutative ring $R$.

    (i) The set $\mathfrak{L} (A,B;C)$ of all $R$-bilinear maps $A\times B \to C$ is an $R$-module with $(f+g)(a,b)=f(a,b)+g(a,b)$ and $(rf)(a,b)=rf(a,b)$.

    (ii) $\Hom_{R}(A\otimes_R B, C)\simeq \Hom_{R}(A, \Hom_{R}(B, C))\simeq \Hom_{R}(B, \Hom_{R}(A, C)) \simeq \mathfrak{L} (A,B;C)$.
\end{theorem}
\begin{proof}
    (ii) Since $R$ is commutative, it is easy to prove $A\otimes_R B \simeq B \otimes_R A$. Then by Theorem3.31, $\Hom_{R}(A\otimes_R B, C) \simeq \Hom_{R}(A, \Hom_{R}(B, C)) \simeq \Hom_{R}(B, \Hom_{R}(A, C))$. The only thing we need to prove is that $\mathfrak{L} (A,B;C)$ is isomorphic to one of them. We now prove $\mathfrak{L} (A,B;C) \simeq \Hom_{R}(A\otimes_R B, C)$.

    For all $f \in \mathfrak{L} (A,B;C)$, by universal property there is a unique $R$-module homomorphism $\bar{f}$ such that $f(a,b)=\bar{f}(a\otimes b)$. Define $\varphi:\mathfrak{L} (A,B;C)\to \Hom_{R}(A \otimes_R B,C )$ such that $\varphi(f)=\bar{f}$, then clearly $\varphi$ is well-defined. $\varphi(f_{1}+f_{2})=\bar{f_{1}+f_{2}}=\bar{f_{1}}+\bar{f_{2}}=\varphi(f_{1})+\varphi(f_{2})$ by the uniqueness of the induced homomorphim. And since $\varphi(rf)(a \otimes b)=(\bar{rf})(a \otimes b)=(rf)(a,b)=r(f(a,b))=r(\bar{f}(a\otimes b))=r\varphi(f)(a \otimes b)$, $\varphi$ is an $R$-module homomorphism. 
    
    Define $\psi:\Hom_{R}(A\otimes_R B, C) \to \mathfrak{L}(A,B;C)$ such that $(\psi(g))(a,b)=g(a \otimes b)$. For $a, a_{i} \in A, b, b_{i} \in B, r \in R$, $(\psi(f))(a_{1}+a_{2},b)=f((a_{1}+a_{2})\otimes b)=f(a_{1}\otimes b)+f(a_{2} \otimes b)=(\psi(f))(a_{1},b)+(\psi(f))(a_{2},b)$ and similar linearity holds for the second varible and $(\psi(f))(ar,b)=f(ar \otimes b)=f(a \otimes rb)=(\psi f)(a,rb)$. Then $\psi$ is well-defined.

    It is easy to verify that $\psi $ is the inverse map of $\varphi$, hence $\mathfrak{L} (A,B;C) \simeq \Hom_{R}(A\otimes_R B, C)$. This completes the proof.
\end{proof}

\begin{theorem}
    Let $F$ be a free module over a p.i.d $R$ and $G$ a submodule of $F$. Then $G$ is a free $R$-module and $\operatorname{rank} G \leqslant  \operatorname{rank} F$.
\end{theorem}
\begin{remark}
    Here $F$ need not to be finite generated. The proof can be seen at Theorem6.1 in the book.
\end{remark}

\begin{corollary}
    A unitary module $A$ over a p.i.d is free iff $A$ is projective.
\end{corollary}
Section 6 primarily deals with the decomposition of finitely generated modules over a p.i.d. The material closely parallels the corresponding sections in standard abstract algebra, where the discussion of decomposition proceeds from the general case to torsion modules, and finally to $p$-primary modules. The only minor differences lie in the proof details; moreover, Hungerford does not employ matrices as a mathematical tool here. Therefore, a detailed discussion of this section is omitted for now, and relevant concepts from Section 6 will only be supplemented when required in subsequent chapters.

Personally, I don't feel this section is very well done. It either completely ignores the material learned in previous sections or merely scratches the surface, utilizing it only in minor parts of the proofs. Although it abandons matrix tools, the resulting proofs are neither more profound nor more concise. Instead, they get bogged down in technicalities, to the point of becoming rather complicated. In summary, I am personally not a big fan of this section (just my personal opinion).

\begin{proposition}
    If $A$ and $B$ are cyclic modules over $R$ of nonzero orders $r$ and $s$ respectively. Let $d=(r,s)$ be the g.c.d of $r,s$and $m=[r,s]$ the l.c.m of $r,s$. Then the invariant factors of $A \oplus B$ are $d$ and $m$.
\end{proposition}
\begin{proof}
    If we can prove $0\to \faktor{R}{(m)}\xrightarrow{f}\faktor{R}{(r)}\oplus \faktor{R}{(s)}\xrightarrow{g}\faktor{R}{(d)}\to 0$ is split, then $\faktor{R}{(r)}\oplus \faktor{R}{(s)} \simeq \faktor{R}{(d)} \oplus \faktor{R}{(m)}$ and this implies that the invariant factors of $A\oplus B$ are $d,m$ since $A \simeq \faktor{R}{(r)}$, $B \simeq \faktor{R}{(s)}$. 

    Define $f(x+(m))=(x+(r),x+(s))$ and $g(x+(r),y+(s))=x-y+(d)$, then verify $f$ and $g$ are well-defined. In fact, suppose $r=r'd$ and $s=s'd$, then $(r',s')=1$ and $m=r's'd$. Thus there exist $u,v \in R$ such that $ur'+vs'=1$. Note that $m(1+(r),1+(s))=0$, then $f$ is well-defined. If $(x_{1}+(r),y_{1}+(s))=(x_{2}+(r),y_{2}+(s))$, then $\exists k_{1},k_{2} \in R$ s.t. $x_{1}-x_{2}=k_{1}r, y_{1}-y_{2}=k_{2}s$. Thus $(x_{1}-y_{1}+(d))-(x_{2}-y_{2}+(d))=k_{1}r-k_{2}s+(d)=(k_{1}r'-k_{2}s')d+(d)=0$. Therefore $g$ is well-defined.
    
    And then verify $f$ is a monomorphism and $g$ is an epimorphism. Define $\alpha:\faktor{R}{(d)}\to \faktor{R}{(r)}\oplus \faktor{R}{(s)}$ such that $\alpha(1+(d))=(ur'+(r),-vs'+(s))$. Note that $g \alpha=\operatorname{id}$. Hence the sequence is split.
\end{proof}


\newpage
\section{Linear Algebra}
\begin{theorem}
    Let $R$ be a ring with identity and let $E,F,G$ be free left $R$-modules with finite ordered bases $U=\{u_{1},u_{2}, \ldots,u_{n}\}$, $V=\{v_{1},v_{2}, \ldots,v_{m}\}$, $W=\{w_{1},w_{2}, \ldots,w_{p}\}$ respectively. If $f \in \Hom_{R}(E, F)$ has $n \times m$ matrix $A$ relative to bases $U$ and $V$ and $g \in \Hom_{R}(F, G)$ has $n \times p$ matrix $B$ relative to bases $V$ and $W$, then $gf \in \Hom_{R}(E, G)$ has $n \times p$ matrix $AB$ relative to bases $U$ and $W$.
\end{theorem}
\begin{proof}
    An intuitive but non-rigorous proof:
    $$g(f(\begin{bmatrix} u_1 \\ u_2 \\ \vdots \\ u_n \end{bmatrix}))=g(A\begin{bmatrix} v_1 \\ v_2 \\ \vdots \\ v_m \end{bmatrix})=Ag(\begin{bmatrix} v_1 \\ v_2 \\ \vdots \\ v_m \end{bmatrix})=AB \begin{bmatrix} w_1 \\ w_2 \\ \vdots \\ w_p \end{bmatrix}.$$
\end{proof}
\begin{definition}
    If $S$ and $T$ are any rings, then a function $\theta:S\to T$ is said to be an \textbf{anti-isomorphism} if $\theta$ is an isomorphism of additive group such that $\theta(s_{1}s_{2})=\theta(s_{2})\theta(s_{1})$ for all $s_{i} \in S$.

    If $R $ is a ring, then the \textbf{opposite ring} of $R$, denoted as $R^{op}$, is the ring that has the same set of elements as $R$, the same addition as $R$, and multiplication $\circ$ given by $a \circ b=ba$ where $ba$ is the product in $R$.
\end{definition}

\begin{theorem}
    Let $R$ be a ring with identity and $E$ a free left $R$-module with a finite basis of $n$ elements. Then there is an isomorphism of rings:
    \begin{equation*}
        \Hom_{R}(E, E) \simeq \Mat_n(R^{op}).
    \end{equation*}
\end{theorem}

\begin{definition}
    Let $B_1, B_2, \ldots, B_n$ and $C$ be modules over a commutative ring $R$ with identity. A function $f:B_{1} \times \cdots \times B_{n} \to C$ is said to be \textbf{$R$-multilinear} if for each $i=1,2, \ldots,n $ and all $r,s \in R$ and $b,b' \in B$, there is $f(b_{1}, \ldots,b_{i-1},rb+sb',b_{i+1}, \ldots,b_{n})=rf(b_{1}, \ldots,b_{i-1},b,b_{i+1}, \ldots,b_{n})+sf(b_{1}, \ldots,b_{i-1},b',b_{i+1}, \ldots,b_{n})$.

    If $C=R$, then f is called an \textbf{$n$-linear} or \textbf{$R$-multilinear form}. If $C=R$ and $B_{1}=B_{2}=\cdots=B_{n}=B$, then $f$ is called an \textbf{$R$-multilinear form on $B$}.
\end{definition}

\begin{definition}
    Let $B$ and $C$ be $R$-modules and $f:B^n \to C$ an $R$-multilinear function. multilinear function. Then $f$ is said to be \textbf{symmetric} if
    $$f(b_{\sigma 1}, \dots, b_{\sigma n}) = f(b_1, \dots, b_n) \quad \text{for every permutation } \sigma \in S_n,$$
    and \textbf{skew-symmetric} if
    $$f(b_{\sigma 1}, \dots, b_{\sigma n}) = (\operatorname{sgn} \sigma) f(b_1, \dots, b_n) \quad \text{for every } \sigma \in S_n.$$
    $f$ is said to be \textbf{alternating} if
    $$f(b_1, \dots, b_n) = 0 \quad \text{whenever } b_i = b_j \quad \text{for some } i \neq j.$$
\end{definition}

\begin{theorem}
    If $B$ and $C$ are modules over a commutative ring $R$ with identity, then every alternating $R$-multilinear function $f : B^n \to C$ is skew-symmetric.
\end{theorem}
If $r+r\neq 0$ for all nonzero $r \in R$, then an $n$-linear form $B^n \to R$ is alternating iff it is skew-symmetric. However, if $\operatorname{char} R=2$, it may be skew-symmetric but not alternating.
\begin{example}
    Let $R=B=\mathbb{Z}_{2}$ and $f:B\to R$ as $f(x,y)=xy$. Obviously $f$ is bilinear and skew-symmetric. But $f(1,1)=1\cdot 1 =1\neq 0$, thus $f$ is not alternating.
\end{example}

\begin{theorem}
    If $R$ is a commutative ring with identity and $r \in R$, then there exists a unique alternating $R$-multilinear form $f : (R^n)^n \to R$ such that $f(\varepsilon_1, \varepsilon_2, \dots, \varepsilon_n) = r$, where $\{\varepsilon_1, \dots, \varepsilon_n\}$ is the standard basis of $R^n$.
\end{theorem}

\begin{definition}
    Let $R$ be a commutative ring with identity. The unique alternating $R$-multilinear form $d : \operatorname{Mat}_n R \to R$ such that $d(I_n) = 1_R$ is called the \textbf{determinant function} on $\operatorname{Mat}_n R$. The \textbf{determinant} of a matrix $A \in \operatorname{Mat}_n R$ is the element $d(A) \in R$ and is denoted $|A|$.
\end{definition}

\begin{theorem}
    Let $\phi : E \to E$ be a linear transformation of an $n$-dimensional vector space $E$ over a field $K$.

    (i) $E$ has a basis relative to which the matrix of $\phi$ is the direct sum of the companion matrices of the invariant factors $q_1, \dots, q_t \in K[x]$ of $\phi$.

    (ii) $E$ has a basis relative to which the matrix of $\phi$ is the direct sum of the companion matrices of the elementary divisors $p_1^{m_{11}}, \dots, p_s^{m_{sk_s}} \in K[x]$ of $\phi$.

    (iii) If the minimal polynomial $q$ of $\phi$ factors as $q = (x - b_1)^{r_1}(x - b_2)^{r_2} \dots (x - b_d)^{r_d}$ ($b_i \in K$), which is always the case if $K$ is algebraically closed, then every elementary divisor of $\phi$ is of the form $(x - b_i)^j$ ($j \le r_i$) and $E$ has a basis relative to which the matrix of $\phi$ is the direct sum of the elementary Jordan matrices associated with the elementary divisors of $\phi$.
\end{theorem}

\begin{theorem}
    Let $\phi : E \to E$ be a linear transformation of an $n$-dimensional vector space over a field $K$ with characteristic polynomial $p_{\phi} \in K[x]$, minimal polynomial $q_{\phi} \in K[x]$, and invariant factors $q_1, \dots, q_t \in K[x]$.

    (i) The characteristic polynomial is the product of the invariant factors; that is, $p_{\phi} = q_1 q_2 \dots q_t = q_1 q_2 \dots q_{t-1} q_{\phi}$.

    (ii) (Cayley-Hamilton) $\phi$ is root of its characteristic polynomial; that is, $p_{\phi}(\phi) = 0$.

    (iii) An irreducible polynomial in K[x] divides $p_{\phi} $ if and only if it divides $q_{\phi}$.

    Conclusions (i)--(iii) are valid, \textit{mutatis mutandis}, for any matrix $A \in \operatorname{Mat}_n K$.
\end{theorem}

\begin{definition}
    Let $\phi : E \to E$ be a linear transformation of a vector space $E$ over a field $K$. A nonzero vector $u \in E$ is an \textbf{eigenvector} (or \textbf{characteristic vector} or \textbf{proper vector}) of $\phi$ if $\phi(u) = ku$ for some $k \in K$. An element $k \in K$ is an \textbf{eigenvalue} (or \textbf{proper value} or \textbf{characteristic value}) of $\phi$ if $\phi(u) = ku$ for some nonzero $u \in E$.
\end{definition}

\begin{theorem}
    Let $\phi : E \to E$ be a linear transformation of a finite dimensional vector space $E$ over a field $K$. Then $\phi$ has a diagonal matrix $D$ relative to some ordered basis of $E$ if and only if the eigenvectors of $\phi$ span $E$. In this case the diagonal entries of $D$ are the eigenvalues of $\phi$ and each eigenvalue $k \in K$ appears on the diagonal $\dim_K C(\phi, k)$ times.
\end{theorem}
\newpage

\section{Category}
\begin{definition}
    A \textbf{category} is a class $\mathscr{C}$ of objects $A,B,C, \ldots$ together with a class of disjoint sets $\Hom_{}(A, B)$, one for each pair of objects in $\mathscr{C}$, and with triple $(A,B,C)$ of objects of $\mathscr{C}$ a function $\Hom_{}(B, C) \times \Hom_{}(A, B) \to \Hom_{}(A, C)$ which satisfy two axioms:

    (i) Associativity. If $f:A \to B, g:B\to C, h:C\to D$ are morphisms of $\mathscr{C}$, then $h \circ (g \circ f)=(h \circ g) \circ f$.
    
    (ii) Identity. For each object $B$ of $\mathscr{C}$ there exists a morphism $1_B:B\to B$ such that for any $f:A\to B$ and $g:B \to C$, $1_B \circ f=f$ and $g \circ 1_B=g$.
\end{definition}

\begin{definition}
    A \textbf{concrete category} is a category $\mathscr{C}$ together with a function $\sigma$ that assigns to each object $A$ of $\mathscr{C}$ a set $\sigma(A)$ (called the underlying set of $A$) such that:

    (i) every morphism $A\to B$ is a function on the underlying sets $\sigma(A)\to \sigma(B)$;

    (ii) the identity morphism of each object $A$ is the identity function on the underlying set $\sigma(A)$;

    (iii) compoistion of morphisms in $\mathscr{C}$ agrees with compostion of functions on the underlying sets.
\end{definition}

\begin{example}
    A multiplicative group $G$ can be considered as a category with only one object $G$. The morhpisms $\Hom_{}(G, G)$ of the category are the elements of $G$. Composition of morphisms $a,b$ is simply the composition $ab$ given by the binary operation in $G$. It is easy to verify this is a category. But this is not a concrete category since the object $G$ has no underlying set.
\end{example}

\begin{definition}
    An object $I$ in a category $\mathscr{C}$ is said to be \textbf{universal} (or \textbf{initial}) if for each object $C$ of $\mathscr{C}$, there exists one and only one morphism $I\to C$. An object $T$ of $\mathscr{C}$ is said to be \textbf{couniversal} (or \textbf{terminal}) if for each object $D$ of $\mathscr{C}$, there exists one and only one morphism $D\to T$.
\end{definition}

\begin{theorem}
    Any two universal (couniversal) objects in a category $\mathscr{C}$ are equivalent.
\end{theorem}

\begin{definition}
    Let $\mathscr{C}$ be a category and $\{A_{i}: i \in I\}$ a family of objects of $\mathscr{C}$. A \textbf{product} for the family $\{A_{i}: i \in I\}$ is an object $P$ of $\mathscr{C}$ together with a family of morphisms $\{\pi_i : P\to A_{i}\}$ such that for each object $B$ of $\mathscr{C}$ and family of morphisms $\{\varphi_i:B \to A_{i}\}$, there exists a unique morphism $\varphi: B\to P$ such that $\pi_i \circ \varphi =\varphi_i$ for all $i \in I$. 
\end{definition}

\begin{definition}
    Let $\mathscr{C}$ be a category and $\{A_{i}: i \in I\}$ a family of objects of $\mathscr{C}$. A \textbf{coproduct} for the amily $\{A_{i}: i \in I\}$ is an object $S$ of $\mathscr{C}$ together with a family of morphisms $\{\iota_i : A_{i}\to  S\}$ such that for each object $B$ of $\mathscr{C}$ and family of morphisms $\{\psi_i: A_{i} \to B\}$, there exists a unique morphism $\psi: B\to P$ such that $\psi \circ \iota_i = \psi_i$ for all $i \in I$. 
\end{definition}

It is obvious that a product of a family of objects in a category is a couniversal or terminal object and a coprouct of a family of objects in a category is a universal or initial object.

\begin{theorem}
    If $(A,\{\varphi_i\})$ and $(B,\{\psi_i\})$ are both products [resp. coproducts] for a family of objects of a category, then $A$ and $B$ are equivalent.
\end{theorem}

\begin{definition}
    Let $F$ be an object in a concrete category $\mathscr{C}$, $X$ a nonempty set, and $i:X\to F$ a map (of sets). F is \textbf{free on the set $X$} provided that for any object $A$ of $\mathscr{C}$ and map (of sets) $f:X\to A$, there exists a unique morphism $\bar{f}:F\to A$ of $\mathscr{C}$ such that $\bar{f} \circ i = f$.
\end{definition}

Note that a free object of a category is also a universal object. By the Theorem 5.1, free objects and coproducts are actually equivalent in a concrete category.

\begin{definition}
    Let $\mathscr{C}$ and $\mathscr{D}$ be categories. A \textbf{contravariant functor} $S$ from $\mathscr{C}$ to $\mathscr{D}$ (denoted $S:\mathscr{C}\to \mathscr{D}$) is a pair functions (both denoted by $S$), an object function that assigns to each object $C$ of $\mathscr{C}$ an object $S(C)$ of $\mathscr{D}$ and a morphism function which assigns each morphism $f:C\to C'$ of $\mathscr{C}$ a morphism $S(f):S(C')\to S(C)$ such that:

    (i) $S(1_C)=1_{S(C)}$ for every indetity morphism $1_C$ of $\mathscr{C}$;

    (ii) $S(g \circ f)=S(f) \circ S(g)$ for any two morphisms $f,g$ of $\mathscr{C}$ whose composition $g \circ f$ is defined.
\end{definition}


\begin{example}
    Let $R$ be a ring and $A$ a fixed $R$-module. For each $R$-module $C$, let $S(C)=\Hom_{R}(C, A)$. For each $R$-module homomorphism $f:C \to C'$, let $S(f)$ be the usual induced map $\bar{f}:\Hom_{R}(C', A)\to \Hom_{R}(C, A)$. Then $S$ is a contravariant functor from the category of $R$-module to the category of abelian group since $\Hom_{R}(C, A)$ is an abelian group.

    More generally, let $A$ be a fixed object in a category $\mathscr{C}$. Define a contravariant functor $h^A=\Hom_{}(-, A)$ from $\mathscr{C}$ to the category $\mathscr{S}$ of sets by assigning to an object $C$ of $\mathscr{C}$ the set $h^A(C)=\Hom_{}(C, A)$ of all morphisms in $\mathscr{C}$ from $C$ to $A$. If $f:C\to C'$ is a morhpism of $\mathscr{C}$, the $h^A(C):\Hom_{}(C', A)\to \Hom_{}(C, A)$ be the function given by $g \mapsto g \circ f$ ($g \in \Hom_{}(C', A)$). The functor $h^A$ is called \textbf{contravariant hom functor}.
\end{example}

The product category of multiple categories is defined in an intuitive manner analogous to the Cartesian product. On a product category, one can define functors of several variables (mapping to a single codomain category). Such a functor can be covariant in certain variables and contravariant in others, with its action on morphisms dictated by the variance of each respective variable. Moreover, if all but one variable are fixed, the multivariable functor naturally reduces to an ordinary single-variable functor.

\begin{example}
    $\Hom_{R}(-, -)$ is a functor of two variables, contravariant in the first and covariant in the second, from the category of $\mathscr{M} \times \mathscr{M}$ of left $R$-module to the category of abelian groups.

    More generally, let $\mathscr{C}$ be any category. Consider the functor that assigns to each pair $(A,B)$ of objects of $\mathscr{C}$ the set $\Hom_{\mathscr{C}}(A, B)$ and to each pair of morphisms $f:A\to A'$ and $g:B\to B'$ the function 
    \begin{equation*}
        \Hom_{}(f, g): \Hom_{\mathscr{C}}(A', B) \to \Hom_{\mathscr{C}}(A, B')
    \end{equation*}
    given by $h \mapsto g \circ h \circ f$. Then $\Hom_{\mathscr{C}}(-, -)$ is a functor of two variables from $\mathscr{C}$ to the category $\mathscr{S}$ of sets, contravariant in the first variable and covariant in the second.
\end{example}

\begin{example}
    Let $K$ be commutative ring with identity. Then the functor given by
    \begin{align*}
        T(A_1, \cdots, A_n) &= A_1 \otimes_K \cdots \otimes_K A_n \\
        T(f_1, \cdots, f_n) &= f_1 \otimes \cdots \otimes f_n
    \end{align*}
    is a functor of $n$ covariant variables from the category of $K$-modules to itself. $T$ is also written as $- \otimes_K \cdots \otimes_K -$.
\end{example}

\begin{definition}
    Let $\mathscr{C}$ and $\mathscr{D}$ be categories and $T:\mathscr{C}\to \mathscr{D}$, $S:\mathscr{C}\to \mathscr{D}$ contravariant functors. A \textbf{natural transformation} $\alpha:S\to T$ is a function that assigns to each object $C$ of $\mathscr{C}$ a morphism $\alpha_C:S(C)\to T(C)$ of $\mathscr{D}$ in such a way that for every morphism $f:C\to C'$ the diagram
    $$    \begin{tikzcd}
        S(C) \arrow[r, "\alpha_C"]                        & T(C)                     \\
        S(C') \arrow[r, "\alpha_{C'}"'] \arrow[u, "S(f)"] & T(C') \arrow[u, "T(f)"']
    \end{tikzcd}$$
    in $\mathscr{D}$ is commutative. If $\alpha_C$ is an equivalence for every $C$ in $\mathscr{C}$, then $\alpha$ is a \textbf{natural isomorphism} of the functors $S$ and $T$.    
\end{definition}

\begin{example}
    Let $T:\mathscr{C}\to \mathscr{C}$ is any functor, the the assignment $C \mapsto 1_{T(C)}$ defines a natural isomorphism $I_T: T\to T$, called the identity natural isomorphism.
\end{example}

\begin{example}
    Let $R$ be a ring with identity and $\mathscr{M}$ the category of left $R$-modules. Then there is a natural isomorphism from the functor $R \otimes -$ to $I_{\mathscr{M}}$ since for each module homomorphism $f:A\to B$, the diagram 
    \begin{equation*}
        \begin{tikzcd}
        R \otimes_R B \arrow[r, "\alpha_B"] \arrow[d, "1_R \otimes f"'] & B \arrow[d, "f"] \\
        R \otimes C \arrow[r, "\alpha_C"']                             & C
        \end{tikzcd}
    \end{equation*}
    is commutative.
\end{example}

\begin{definition}
    Let $S$ be a contravariant functor from a category $\mathscr{C}$ to the category $\mathscr{S}$ of sets. $S$ is said to be a \textbf{representable functor} if there is an object $A$ of $\mathscr{C}$ and a natural isomorphism $\alpha$ from the contravariant hom functor $h^A=\Hom_{\mathscr{C}}(-,A)$ to $S(-)$. The pair $(A,\alpha)$ is called a representation of $S$ and $S$ is said to be represented by the object $A$.
\end{definition}

\begin{example}
    Let $\mathscr{G}$ be the category of groups and $F$ the forgetful functor from $\mathscr{G}$ to the category $\mathscr{S}$ of the set (that is, $F$ assigns to a group its underlying set). We need find a group $A$ of $\mathscr{G}$ so that there is a natural isomorphism from $\Hom_{\mathscr{G}}(A,-)$ to $F(-)$. Take $A=\mathbb{Z}$, then for each $\alpha_G \in \Hom_{\mathscr{G}}(\mathbb{Z},G)$ consider $\alpha_{G,g}(1)=g$ and $\alpha_{G,g}(n)=g^n$ where $G \in \mathscr{G}$ and $g \in G$. By the property of the homomorphism from $\mathbb{Z}$ to a group, every $\alpha_G$ is determined by the value $g=\alpha_{G,g}(1)$, that is, the homomorphisms in $\Hom_{\mathscr{G}}(\mathbb{Z},G)$ one to one correspond with the elements in $F(G)$ and then $\Hom_{\mathscr{G}}(\mathbb{Z},G)\simeq F(G)$. This implies that the forgetful functor $F$ is representable and every element in a group $G$ corresponds to a morphism from $\mathbb{Z}$ to $G$. 
    
    Here is a question. We do make a great transform from specific elements in a group to the group homomorphisms from $\mathbb{Z}$ to the group $G$. What about the operation in $G$? Is there a group $A$ such that there is a natural isomorphism from $\Hom_{\mathscr{G}}(A,-)$ to $\operatorname{id}_{\mathscr{G}}(-)$?
\end{example}

\begin{example}
    
\end{example}

\begin{definition}
    Let $T:\mathscr{C}\to \mathscr{D}$ be a covariant functor and $\mathscr{C}_T$ the category with objects all pairs $(C,s)$ where $C$ is an object of $\mathscr{C}$ and $s \in T(C)$. A morphism in $\mathscr{C}_T$ from $(C,s)$ to $(D,t)$ is defined to be a morphism $f:C\to D$ of $\mathscr{C}$ such that $T(f)(s)=t \in T(D)$. A universal object in $\mathscr{C}_T$ is called a \textbf{universal element} of the functor $T$.
\end{definition}

\begin{example}
    Consider the example of Example 5.7. Let $\mathscr{G}_F$ be the category with objects all pairs $(G,g)$ where $G$ is an object of $\mathscr{G}$ and $g \in F(G)$. Consider the pair $(\mathbb{Z},1)$, then for each $(G,g)$ of $\mathscr{G}_F$ there is a morphism $\varphi:(\mathbb{Z},1)\to (G,g)$ such that $\varphi(1)=g$. Hence $(\mathbb{Z},1)$ is the universal object in $\mathscr{G}_F$, and then a universal element of the forgetful functor $F$.
\end{example}

\begin{lemma}
    Let $T:\mathscr{C}\to \mathscr{S}$ be a contravariant functor from a category $\mathscr{C}$ to the category $\mathscr{S}$ of sets and let $A$ be an object of $\mathscr{C}$.

    (i) If $\alpha:h^A\to T$ is a natural transformation for the contravariant functor $h^A=\Hom_{\mathscr{C}}(-,A)$ to $T$ and $u=\alpha_A(1_A) \in T(A)$, then for any object $C$ of $\mathscr{C}$ and $g \in \Hom_{\mathscr{C}}(C,A)$, there is $$\alpha_C(g)=T(g)(u).$$

    (ii) If $u \in T(A)$ and for each object $C$ of $\mathscr{C}$, $\beta_C:\Hom_{\mathscr{C}}(C,A)\to T(C)$ is the map defined by $\beta_C(g)= T(g)(u)$, then there is a natural transformation $\beta:h^A\to T$ such that $\beta_A(1_A)=u$.
\end{lemma}

\begin{theorem}
    Let $T:\mathscr{C}\to \mathscr{S}$ be a covariant [resp. contravariant] functor from a category $\mathscr{C}$ to the category $\mathscr{S}$ of sets. Then there is a one-to-one correspondence between the class $X$ of all representations of $T$ and the class $Y$ of all universal elements of $T$, given by $(A,\alpha)\mapsto (A,\alpha_A(1_A))$. 
\end{theorem}

\begin{corollary}(Yoneda)
    Let $T:\mathscr{C}\to \mathscr{S}$ be a covariant [resp. contravariant] functor from a category $\mathscr{C}$ to the category $\mathscr{S}$ of sets and let $A$ be an object of $\mathscr{C}$. Then there is a one-to-one correspondence between the set $T(A)$ and the set $\operatorname{Nat}_{}(h_A,T)$ of all natural transformations from the covariant functor $h_A=\Hom_{\mathscr{C}}(A,-)$ to the functor $T$. This bijection is natural in $A$ and $T$.
\end{corollary}

\begin{theorem}
    Let $\mathscr{C}$ and $\mathscr{D}$ be categories and $T$ a functor from the product category $\mathscr{C} \times \mathscr{D}$ to the category $\mathscr{S}$ of sets, contravariant in the first variable and covariant in the second, such that for each object $C$ of $\mathscr{C}$, the covariant functor $T(C,-) : \mathscr{D} \to \mathscr{S}$ has a representation $(A_C, \alpha^C)$. Then there is a unique covariant functor $S : \mathscr{C} \to \mathscr{D}$ such that $S(C) = A_C$ and there is a natural isomorphism from $\Hom_{\mathscr{D}}(S(-), -)$ to $T$, given by
    $$\alpha^C_D : \operatorname{hom}_{\mathscr{D}}(S(C), D) \to T(C,D).$$
\end{theorem}
\begin{proof}
    Define $S(C)=A_C$ then we need to define $S(f)$ for each morphism $f:C\to C'$ of $\mathscr{C}$. Note that 
    $$\begin{aligned}
    &T(f,-): T(C',-)\to T(C,-), \\
    &\alpha^C: h_{A_{C}}=\Hom_{\mathscr{D}}(A_C,-) \simeq T(C,-), \\
    &\alpha^{C'}: h_{A_{C'}}=\Hom_{\mathscr{D}}(A_{C'},-) \simeq T(C',-).
    \end{aligned}$$
    Then let the functor $U$ be 
    $$U=(\alpha^C)^{-1} \circ T(f,-) \circ \alpha^{C'}:\Hom_{\mathscr{D}}(A_{C'},-)\to \Hom_{\mathscr{D}}(A_C,-),$$
    hence $U \in \Nat(h_{A_{C'}},h_{A_C})$.

    By Yoneda's theorem there is a one-to-one correspondence between $\Nat(h_{A_{C'}},h_{A_C})$ and $h_{A_C}(A_{C'})$. Thus there exists $h \in h_{A_C}(A_{C'})=\Hom_{\mathscr{D}}(A_C,A_{C'})$ corresponding with the functor $U$. In fact, $h=U_{A_{C'}}(1_{A_{C'}})$. Then we define $S(f)=h$. It is easy to verify $S$ is a functor from $\mathscr{C}$ to $\mathscr{D}$ since $S(1_C)=1_{A_C}$ holds by the property of the universal elements and $S(f_{1} \circ f_{2})=S(f_{1}) \circ S(f_{2})$ holds by the construction of $S$.
\end{proof}
\begin{remark}
    Note that for any $f:C\to C'$, 
    $$\begin{aligned}
    S(f) = (\alpha^C_{A_{C'}})^{-1} \Big( T(f, A_{C'}) \big( \alpha^{C'}_{A_{C'}}(1_{A_{C'}}) \big) \Big), \\
    \alpha^C_{A_{C'}}\big(S(f)\big) = T(f, A_{C'}) \big( \alpha^{C'}_{A_{C'}}(1_{A_{C'}}) \big).
    \end{aligned}$$
    Then for $C \xrightarrow{f_2} C' \xrightarrow{f_1} C''$,
    $$\alpha^C_{A_{C''}}\big(S(f_1 \circ f_2)\big) = T(f_1 \circ f_2, A_{C''}) \big( \alpha^{C''}_{A_{C''}}(1_{A_{C''}}) \big) .$$
    Note that
    $$\begin{aligned}
    \alpha^C_{A_{C''}}\big(S(f_1) \circ S(f_2)\big) 
    &= T(1_C, S(f_1)) \Big[ \alpha^C_{A_{C'}}\big(S(f_2)\big) \Big]
    = T(1_C, S(f_1)) \Big[ T(f_2, A_{C'}) \big( \alpha^{C'}_{A_{C'}}(1_{A_{C'}}) \big) \Big] \\
    &= T(f_2, A_{C''}) \Big[ T(1_{C'}, S(f_1)) \big( \alpha^{C'}_{A_{C'}}(1_{A_{C'}}) \big) \Big]
    = T(f_1 \circ f_2, A_{C''}) \big( \alpha^{C''}_{A_{C''}}(1_{A_{C''}}) \big).
    \end{aligned}$$
    Hence $S(f_1 \circ f_2) = S(f_1) \circ S(f_2)$.

    Too complicated!!!
\end{remark}









\end{document}

















